% ============================================================================
%  Plataforma Unooria — Estudo PU-EST-001
%  Teste de carga para uma rede de 800 mil usuários: o que medir, com que
%  números e como executar.
%
%  Compilação (a partir desta pasta):
%     tectonic relatorio-teste-de-carga.tex
%  Reaproveita o preâmbulo da PU-DOC-001 (../../preambulo) para ter as mesmas
%  cores, caixas, tabelas e a bibliografia ABNT.
% ============================================================================
\documentclass[11pt, a4paper, oneside, openany]{report}

% ============================================================================
%  estilo.tex — identidade visual, tipografia e componentes do documento
%
%  Tipografia: Roboto (texto e titulos) + Roboto Mono (codigo). As duas vem do
%  pacote `roboto`, entao o documento compila com pdfLaTeX, XeLaTeX ou Tectonic
%  sem fontspec explicito.
%
%  Paleta: os tokens `--unooria-*` do @unooria/design-system, os MESMOS valores
%  que `packages/ui/styles/plataforma.css` usa como fallback. Cor nova neste
%  documento so entra se existir como token no pacote.
% ============================================================================

% --- Codificacao e idioma ---------------------------------------------------
\usepackage[T1]{fontenc}
\usepackage[utf8]{inputenc}
\usepackage[shorthands=off,brazil]{babel}   % shorthands OFF: com babel-pt o `"` fica ativo e quebra codigo em linha
\usepackage{csquotes}

% --- Tipografia -------------------------------------------------------------
\usepackage[sfdefault]{roboto}            % Roboto como familia padrao
\usepackage[scaled=0.90]{roboto-mono}     % Roboto Mono para codigo
\usepackage{microtype}
\usepackage{textcomp}

% --- Pagina -----------------------------------------------------------------
\usepackage[
  a4paper,
  top=2.6cm, bottom=2.5cm, left=2.6cm, right=2.4cm,
  headheight=15pt, headsep=18pt, footskip=26pt
]{geometry}
\usepackage{pdflscape}
\usepackage{setspace}
\onehalfspacing

% --- Cores canônicas do Brandbook Unooria ("BLUNOORIA") -----------------------
% Regra semântica: "A cor não decora. Ela organiza."
%   - Azul Profundo (#0E1039): Estrutura, corpo, regras e títulos
%   - Lavanda (#8B85AC): Inteligência, decisões arquiteturais (ADR), IA
%   - Azul Claro (#7FA5E2): Interface, navegação, referências e notas
\usepackage{xcolor}
\definecolor{unooNavy}        {HTML}{0E1039}  % Azul Profundo — Estrutura (C100 M96 Y44 K56)
\definecolor{unooLavender}    {HTML}{8B85AC}  % Lavanda — Inteligência (C52 M48 Y15 K2)
\definecolor{unooBlue}        {HTML}{7FA5E2}  % Azul Claro — Interface (C54 M29 Y0 K0)
\definecolor{unooBlueLight}   {HTML}{B8D0F5}  % Variação clara de interface
\definecolor{unooBlueWash}    {HTML}{F0F4FC}  % Fundo suave de interface
\definecolor{unooLavenderWash}{HTML}{F3F2F8}  % Fundo suave de inteligência
\definecolor{unooLavenderSoft}{HTML}{5D577D}  % Lavanda escurecido p/ alto contraste tipográfico
\definecolor{unooInk}         {HTML}{0E1039}  % neutral-950 / texto primário ancorado na marca
\definecolor{unooInkSoft}     {HTML}{3A3C59}  % neutral-700 / texto secundário
\definecolor{unooMuted}       {HTML}{6D6E82}  % neutral-500 / legendas e metadados
\definecolor{unooGrid}        {HTML}{E3E5EE}  % Linhas de grade e bordas sutis
\definecolor{unooAxis}        {HTML}{CACDDC}  % Réguas e separadores
\definecolor{unooSurface}     {HTML}{FFFFFF}  % Fundo puro
\definecolor{unooPaper}       {HTML}{F7F8FC}  % Fundo alternado / zebra
% Tons funcionais (alertas de sistema e status):
\definecolor{unooRed}         {HTML}{C92A2A}  % danger-600
\definecolor{unooRedWash}     {HTML}{FDF2F2}
\definecolor{unooOrange}      {HTML}{D97706}  % warning-600
\definecolor{unooOrangeWash}  {HTML}{FFFBEB}
\definecolor{unooGreen}       {HTML}{2F855A}  % Status funcional aprovado
\definecolor{unooGreenWash}   {HTML}{F0FDF4}
\definecolor{unooYellow}      {HTML}{B45309}
\definecolor{unooYellowWash}  {HTML}{FEF3C7}
\definecolor{unooPink}        {HTML}{9D4EDD}

% --- Tabelas ----------------------------------------------------------------
\usepackage{booktabs}
\usepackage{tabularx}
\usepackage{longtable}
\usepackage{array}
\usepackage{multirow}
\usepackage{makecell}
\usepackage{colortbl}
\usepackage{ltablex}          % tabularx com quebra de pagina
\keepXColumns
\renewcommand{\arraystretch}{1.22}
\arrayrulecolor{unooAxis}
\setlength{\aboverulesep}{0pt}
\setlength{\belowrulesep}{0pt}
\setlength{\extrarowheight}{2pt}

% Tabelas em \small por decisao de composicao: as celulas deste documento
% misturam prosa e identificadores monoespacados longos (caminhos de arquivo,
% nomes de coluna), que NAO hifenizam. Em corpo 11 elas transbordavam a caixa;
% em \small cabem, e a tabela continua legivel porque a densidade e alta de
% proposito. `\emergencystretch` cobre o residuo.
\usepackage{etoolbox}
\AtBeginEnvironment{tabularx}{\small\setlength{\emergencystretch}{1.2em}}
\AtBeginEnvironment{longtable}{\small\setlength{\emergencystretch}{1.2em}}

\newcolumntype{L}[1]{>{\raggedright\arraybackslash}p{#1}}
\newcolumntype{C}[1]{>{\centering\arraybackslash}p{#1}}
\newcolumntype{R}[1]{>{\raggedleft\arraybackslash}p{#1}}
\newcolumntype{Y}{>{\raggedright\arraybackslash}X}

% Colunas X sem justificacao: as celulas misturam prosa e codigo monoespacado,
% que nao hifeniza — justificar produziria espacos internos muito largos.
\renewcommand{\tabularxcolumn}[1]{>{\raggedright\arraybackslash}p{#1}}

\newcommand{\thd}[1]{\textcolor{white}{\textbf{#1}}}
\newcommand{\headrow}{\rowcolor{unooNavy}}
\newcommand{\zebra}{\rowcolor{unooPaper}}

% --- Listas -----------------------------------------------------------------
\usepackage{enumitem}
\setlist[itemize]{leftmargin=1.35em, itemsep=2pt, topsep=4pt, parsep=0pt,
                  label=\textcolor{unooBlue}{\raisebox{0.18ex}{\rule{0.40em}{0.40em}}}}
\setlist[enumerate]{leftmargin=1.7em, itemsep=2pt, topsep=4pt, parsep=0pt,
                    font=\bfseries\color{unooNavy}}
\setlist[description]{leftmargin=1.4em, style=nextline,
                      font=\bfseries\color{unooNavy}}

% --- Matematica e unidades --------------------------------------------------
\usepackage{amsmath}
\usepackage{amssymb}
\usepackage{siunitx}
\sisetup{output-decimal-marker={,}, group-separator={.}, group-minimum-digits=4,
         per-mode=symbol, detect-weight=true, detect-family=true}
\DeclareSIUnit{\tokens}{tokens}
\DeclareSIUnit{\req}{req}

% --- Figuras e legendas -----------------------------------------------------
\usepackage{graphicx}
\usepackage{float}
\usepackage{adjustbox}
\usepackage[font=small, labelfont={bf,color=unooNavy}, textfont=normalfont,
            labelsep=endash, justification=raggedright, singlelinecheck=false,
            skip=8pt]{caption}
\usepackage{subcaption}
\captionsetup[table]{position=top, skip=6pt}

% Ajusta a figura a largura do texto apenas quando ela excede a caixa.
\newcommand{\fitwidth}[1]{\adjustbox{max width=\linewidth, center}{#1}}
\newcommand{\fitpage}[1]{\adjustbox{max width=\linewidth, max height=0.82\textheight, center}{#1}}

% --- Blocos de destaque (tcolorbox) ----------------------------------------
\usepackage[most]{tcolorbox}
\tcbuselibrary{skins,breakable}

\tcbset{
  unoobase/.style={
    enhanced, breakable, sharp corners,
    boxrule=0pt, leftrule=2.6pt, arc=0pt,
    left=10pt, right=10pt, top=7pt, bottom=7pt,
    fonttitle=\bfseries\small,
    before skip=10pt, after skip=10pt
  }
}

% `colbacktitle` e OBRIGATORIO aqui: sem ele o tcolorbox pinta a faixa do
% titulo com a cor de `colframe`, e o titulo (que usa a MESMA cor em
% `coltitle`) fica invisivel — texto vermelho sobre faixa vermelha. O sintoma
% e uma barra colorida solida no lugar do rotulo do bloco.
\newtcolorbox{nota}[1][]{unoobase, colback=unooBlueWash, colframe=unooBlue,
  colbacktitle=unooBlueWash, coltitle=unooNavy, titlerule=0pt, title={NOTA}, #1}
\newtcolorbox{atencao}[1][]{unoobase, colback=unooOrangeWash, colframe=unooOrange,
  colbacktitle=unooOrangeWash, coltitle=unooOrange, titlerule=0pt,
  title={ATENÇÃO}, #1}
\newtcolorbox{decisao}[1][]{unoobase, colback=unooLavenderWash, colframe=unooLavender,
  colbacktitle=unooLavenderWash, coltitle=unooNavy, titlerule=0pt,
  title={DECISÃO DE PROJETO}, #1}
\newtcolorbox{limitacao}[1][]{unoobase, colback=unooRedWash, colframe=unooRed,
  colbacktitle=unooRedWash, coltitle=unooRed, titlerule=0pt,
  title={LIMITAÇÃO CONHECIDA}, #1}
\newtcolorbox{regra}[1][]{unoobase, colback=unooRedWash, colframe=unooRed,
  colbacktitle=unooRedWash, coltitle=unooRed, titlerule=0pt,
  title={REGRA INVIOLÁVEL}, #1}
\newtcolorbox{destaque}[1][]{unoobase, colback=unooPaper, colframe=unooLavender, #1}
\newtcolorbox{especificacao}[1]{unoobase, colback=unooSurface, colframe=unooNavy,
  colbacktitle=unooSurface, coltitle=unooNavy, titlerule=0pt, title={#1}}

% --- Codigo-fonte -----------------------------------------------------------
\usepackage{listings}

\definecolor{lstBg}     {HTML}{FAFBFB}
\definecolor{lstFrame}  {HTML}{E4E9EB}
\definecolor{lstKeyword}{HTML}{2F4F7F}
\definecolor{lstString} {HTML}{0F6B52}
\definecolor{lstComment}{HTML}{617379}
\definecolor{lstEmph}   {HTML}{A9490F}

\lstdefinestyle{unoo}{
  backgroundcolor=\color{lstBg},
  basicstyle=\ttfamily\scriptsize\color{unooInk},
  keywordstyle=\color{lstKeyword}\bfseries,
  keywordstyle=[2]\color{lstEmph}\bfseries,
  keywordstyle=[3]\color{unooBlue},
  stringstyle=\color{lstString},
  commentstyle=\color{lstComment}\itshape,
  numberstyle=\ttfamily\tiny\color{unooMuted},
  identifierstyle=\color{unooInk},
  numbers=left, numbersep=9pt, stepnumber=1,
  frame=single, framesep=6pt, rulecolor=\color{lstFrame},
  xleftmargin=17pt, xrightmargin=2pt,
  breaklines=true, breakatwhitespace=false, breakindent=0pt,
  postbreak=\mbox{\textcolor{unooBlue}{$\hookrightarrow$}\space},
  showstringspaces=false, tabsize=2,
  columns=fullflexible, keepspaces=true, upquote=true,
  aboveskip=10pt, belowskip=10pt,
  literate=
    {á}{{\'a}}1 {à}{{\`a}}1 {ã}{{\~a}}1 {â}{{\^a}}1
    {é}{{\'e}}1 {ê}{{\^e}}1 {í}{{\'i}}1
    {ó}{{\'o}}1 {õ}{{\~o}}1 {ô}{{\^o}}1
    {ú}{{\'u}}1 {ü}{{\"u}}1 {ç}{{\c c}}1
    {Á}{{\'A}}1 {À}{{\`A}}1 {Ã}{{\~A}}1 {Â}{{\^A}}1
    {É}{{\'E}}1 {Ê}{{\^E}}1 {Í}{{\'I}}1
    {Ó}{{\'O}}1 {Õ}{{\~O}}1 {Ô}{{\^O}}1
    {Ú}{{\'U}}1 {Ç}{{\c C}}1
    {°}{{\textdegree}}1 {º}{{\textordmasculine}}1
    {µ}{{$\mu$}}1 {×}{{$\times$}}1 {≈}{{$\approx$}}1
    {≥}{{$\geq$}}1 {≤}{{$\leq$}}1 {≠}{{$\neq$}}1
    {→}{{$\rightarrow$}}1 {←}{{$\leftarrow$}}1 {⇒}{{$\Rightarrow$}}1
    {—}{{---}}3 {–}{{--}}2 {…}{{\ldots}}3
    {“}{{\textquotedblleft}}1 {”}{{\textquotedblright}}1
    {‘}{{\textquoteleft}}1 {’}{{\textquoteright}}1
    {ª}{{\textordfeminine}}1 {§}{{\S}}1
}
\lstset{style=unoo}

\lstdefinelanguage{tsunoo}[]{Java}{
  morekeywords={const,let,export,default,import,from,as,interface,type,await,
                async,function,return,string,number,boolean,readonly,unknown,
                never,satisfies,enum,declare,namespace,implements},
  morekeywords=[2]{z,Prisma,useEffect,postMessage},
  sensitive=true,
  morecomment=[l]{//},
  morecomment=[s]{/*}{*/},
  morestring=[b]",
  morestring=[b]',
  morestring=[b]`
}
\lstdefinelanguage{sqlunoo}[]{SQL}{
  morekeywords={POLICY,ENABLE,FORCE,ROW,LEVEL,SECURITY,USING,PARTITION,RANGE,
                CONFLICT,DO,NOTHING,JSONB,TIMESTAMPTZ,current_setting,
                set_config,ATTACH,DEFAULT,GENERATED,ALWAYS},
  sensitive=false,
  morecomment=[l]{--},
  morestring=[b]'
}
\lstdefinelanguage{yamlunoo}{
  morekeywords={true,false,null,services,image,build,command,environment,
                volumes,ports,depends_on,healthcheck,condition,runs-on,steps,
                uses,needs,jobs,name,on,with,env,if,run,concurrency,
                permissions,workflow_call,workflow_dispatch},
  sensitive=true,
  morecomment=[l]{\#},
  morestring=[b]",
  morestring=[b]'
}
\lstdefinelanguage{jsonunoo}{
  morekeywords={true,false,null},
  keywordstyle=\color{lstEmph}\bfseries,
  sensitive=true,
  morestring=[b]"
}
\lstdefinelanguage{mermaid}{
  morekeywords={flowchart,graph,sequenceDiagram,erDiagram,classDiagram,gantt,
            subgraph,end,participant,actor,alt,else,loop,par,opt,note,activate,
            deactivate,direction,TB,LR,RL,BT,Person,Person_Ext,System,System_Ext,
            SystemDb,Container,ContainerDb,ContainerQueue,Component,ComponentDb,
            Rel,Rel_D,Rel_U,Rel_L,Rel_R,System_Boundary,Container_Boundary,
            Enterprise_Boundary,Boundary,Deployment_Node,UpdateLayoutConfig,
            UpdateRelStyle,classDef,class,style,state,C4Context,C4Container,
            C4Component,C4Deployment,stateDiagram,stateDiagram-v2},
  sensitive=true,
  morecomment=[l]{\%\%},
  morestring=[b]"
}
\lstdefinelanguage{bashunoo}[]{bash}{
  morekeywords={pnpm,docker,gcloud,node,psql,tectonic,turbo,prisma,gsutil},
  sensitive=true
}

% Codigo em linha.
%
% Uso obrigatorio de \texttt e \textcolor (e nao de {\ttfamily\color{...}}):
% os comandos \text... emitem \leavevmode antes de trocar a fonte, o que faz o
% paragrafo comecar na fonte do texto corrente. Sem isso, uma celula de tabela
% que inicia com codigo teria a primeira linha de base calculada com as metricas
% da fonte monoespacada e ficaria visivelmente desalinhada das demais colunas.
\newcommand{\cd}[1]{\texttt{\textcolor{unooNavy}{#1}}}
\newcommand{\arq}[1]{\texttt{\textcolor{unooInkSoft}{#1}}}
\newcommand{\rota}[1]{\texttt{\textcolor{unooGreen}{#1}}}
\newcommand{\env}[1]{\texttt{\textcolor{lstEmph}{#1}}}
% Ponto de quebra de linha sem hifen, para caminho longo em celula estreita:
% identificador monoespacado nao hifeniza, e sem isto ele transborda a caixa.
\newcommand{\B}{\discretionary{}{}{}}

% --- Titulos ----------------------------------------------------------------
\usepackage{titlesec}
\usepackage{titletoc}

\titleformat{\chapter}[display]
  {\normalfont}
  {\raggedleft\fontsize{62}{62}\selectfont\color{unooBlueLight}\textbf{\thechapter}}
  {-10pt}
  {\raggedright\fontsize{23}{27}\selectfont\bfseries\color{unooNavy}}
  [\vspace{5pt}{\color{unooBlue}\rule{\linewidth}{2.2pt}}]
\titlespacing*{\chapter}{0pt}{-30pt}{26pt}

\titleformat{name=\chapter,numberless}[display]
  {\normalfont}
  {}
  {0pt}
  {\raggedright\fontsize{23}{27}\selectfont\bfseries\color{unooNavy}}
  [\vspace{5pt}{\color{unooBlue}\rule{\linewidth}{2.2pt}}]

\titleformat{\section}
  {\normalfont\Large\bfseries\color{unooNavy}}{\thesection}{0.7em}{}
\titleformat{\subsection}
  {\normalfont\large\bfseries\color{unooInk}}{\thesubsection}{0.7em}{}
\titleformat{\subsubsection}
  {\normalfont\normalsize\bfseries\color{unooInkSoft}}{\thesubsubsection}{0.7em}{}
\titleformat{\paragraph}[runin]
  {\normalfont\normalsize\bfseries\color{unooNavy}}{}{0pt}{}[\;---\;]

\titlespacing*{\section}{0pt}{20pt}{8pt}
\titlespacing*{\subsection}{0pt}{15pt}{6pt}
\titlespacing*{\subsubsection}{0pt}{12pt}{5pt}
\titlespacing*{\paragraph}{0pt}{9pt}{0.4em}

\setcounter{secnumdepth}{3}
\setcounter{tocdepth}{2}

% --- Sumario ----------------------------------------------------------------
\titlecontents{chapter}[0em]
  {\addvspace{9pt}\bfseries\color{unooNavy}}
  {\contentslabel[\thecontentslabel]{1.5em}}
  {}
  {\hfill\contentspage}
\titlecontents{section}[1.7em]
  {\addvspace{1.5pt}}
  {\contentslabel{2.2em}}
  {}
  {\titlerule*[0.7pc]{\color{unooAxis}.}\contentspage}
\titlecontents{subsection}[4.0em]
  {\addvspace{0.5pt}\small\color{unooInkSoft}}
  {\contentslabel{2.8em}}
  {}
  {\titlerule*[0.7pc]{\color{unooGrid}.}\contentspage}

% --- Cabecalho e rodape -----------------------------------------------------
\usepackage{fancyhdr}
\usepackage{lastpage}

\newcommand{\docid}{PU-DOC-001}
\newcommand{\docversao}{1.2}

\renewcommand{\headrule}{{\color{unooGrid}\hrule height \headrulewidth}}
\renewcommand{\footrule}{{\color{unooGrid}\hrule height \footrulewidth}}

\fancypagestyle{unoomain}{%
  \fancyhf{}
  \fancyhead[L]{\footnotesize\color{unooInkSoft}\nouppercase{\leftmark}}
  \fancyhead[R]{\raisebox{-0.15ex}{\includegraphics[height=1.8ex]{figuras/marca/simbolo_azul_profundo.png}}\hspace{4pt}\footnotesize\bfseries\color{unooNavy}Unooria}
  \fancyfoot[L]{\footnotesize\color{unooMuted}\docid{} \textperiodcentered{} v\docversao}
  \fancyfoot[R]{\footnotesize\color{unooInkSoft}\thepage\,/\,\pageref{LastPage}}
  \renewcommand{\headrulewidth}{0.5pt}
  \renewcommand{\footrulewidth}{0.5pt}
}
% Pre-textual: numeracao romana, sem total de paginas (o total do
% \pageref{LastPage} refere-se a numeracao arabe do corpo do documento).
\fancypagestyle{unoofront}{%
  \fancyhf{}
  \fancyhead[L]{\footnotesize\color{unooInkSoft}\nouppercase{\leftmark}}
  \fancyhead[R]{\raisebox{-0.15ex}{\includegraphics[height=1.8ex]{figuras/marca/simbolo_azul_profundo.png}}\hspace{4pt}\footnotesize\bfseries\color{unooNavy}Unooria}
  \fancyfoot[L]{\footnotesize\color{unooMuted}\docid{} \textperiodcentered{} v\docversao}
  \fancyfoot[R]{\footnotesize\color{unooInkSoft}\thepage}
  \renewcommand{\headrulewidth}{0.5pt}
  \renewcommand{\footrulewidth}{0.5pt}
}
\fancypagestyle{plain}{%
  \fancyhf{}
  \fancyfoot[R]{\footnotesize\color{unooInkSoft}\thepage}
  \renewcommand{\headrulewidth}{0pt}
  \renewcommand{\footrulewidth}{0pt}
}
\pagestyle{unoofront}
\renewcommand{\chaptermark}[1]{\markboth{\thechapter\quad #1}{}}

% --- Referencias cruzadas e PDF --------------------------------------------
\usepackage[
  colorlinks=true,
  linkcolor=unooNavy,
  citecolor=unooBlue,
  urlcolor=unooBlue,
  filecolor=unooBlue,
  bookmarksnumbered=true,
  bookmarksopen=true,
  bookmarksopenlevel=1,
  pdfdisplaydoctitle=true
]{hyperref}
\usepackage{bookmark}
\usepackage{url}
\urlstyle{sf}

% --- Bibliografia (ABNT NBR 6023) ------------------------------------------
%
% `abntex2cite` + BibTeX, e nao biblatex+biber: o pipeline de compilacao
% (Tectonic) chama o BibTeX classico automaticamente, sem ferramenta externa.
% A saida sai formatada em ABNT do mesmo jeito.
\usepackage[alf]{abntex2cite}

% --- Semantica de referencia do projeto ------------------------------------
\newcommand{\rf}[1]{\textbf{\textcolor{unooNavy}{#1}}}
\newcommand{\adrr}[1]{\textbf{\textcolor{unooLavenderSoft}{#1}}}
\newcommand{\ucr}[1]{\textbf{\textcolor{unooBlue}{#1}}}
\newcommand{\drr}[1]{\textbf{\textcolor{unooOrange}{#1}}}
\newcommand{\yes}{\textcolor{unooGreen}{\textbf{\checkmark}}}
\newcommand{\no}{\textcolor{unooOrange}{\textbf{--}}}

% Setas: a Roboto NAO tem U+2192 nem U+21D2 (medido — o XeTeX avisa
% "Missing character" e o glifo simplesmente nao e desenhado). Em prosa,
% sempre estes comandos; nunca a seta literal.
\newcommand{\seta}{\ensuremath{\rightarrow}}
\newcommand{\implica}{\ensuremath{\Rightarrow}}
\newcommand{\setabi}{\ensuremath{\leftrightarrow}}

% Espacamento de paragrafo
\setlength{\parindent}{0pt}
\setlength{\parskip}{7pt}

\widowpenalty=10000
\clubpenalty=10000
\flushbottom

\graphicspath{{../../}}
\usepackage{pgfplots}
\pgfplotsset{compat=1.18}
% URLs longas da bibliografia quebram em qualquer caractere, e as páginas não
% esticam o espaço vertical para encher o pé (a bibliografia ficava com
% buracos entre entradas).
\usepackage{xurl}
\raggedbottom

\renewcommand{\docid}{PU-EST-001}
\renewcommand{\docversao}{0.1}

\hypersetup{
  pdftitle={Teste de carga da Plataforma Unooria para uma rede de 800 mil usuarios},
  pdfauthor={Lucas Salomao},
  pdfsubject={Modelo de carga, tipos de teste, criterios de aprovacao e plano de execucao},
  pdfcreator={Tectonic / XeTeX}
}

\begin{document}

% ------------------------------------------------------------------- capa
\begin{titlepage}
\thispagestyle{empty}
\begin{tikzpicture}[remember picture, overlay]
  \fill[unooNavy] (current page.north west) rectangle
        ([yshift=-12mm]current page.north east);
  \fill[unooLavender] ([yshift=-12mm]current page.north west) rectangle
        ([yshift=-13.8mm]current page.north east);
  \fill[unooBlue] ([yshift=-13.8mm]current page.north west) rectangle
        ([yshift=-15.0mm]current page.north east);
\end{tikzpicture}
\vspace*{14mm}

\noindent\includegraphics[width=50mm]{figuras/marca/logo_azul_profundo.png}\par
\vspace{18mm}

{\footnotesize\bfseries\color{unooLavenderSoft} ESTUDO TÉCNICO \quad\textperiodcentered\quad PU-EST-001\par}
\vspace{4mm}
{\fontsize{26}{31}\selectfont\bfseries\color{unooNavy}
 Teste de carga para uma rede\\de 800 mil usuários\par}
\vspace{6mm}
{\large\color{unooInkSoft}
 Como converter usuários cadastrados em carga simultânea, o que medir,
 com quais critérios de aprovação e em que ordem executar\par}

\vfill
\begin{tabularx}{\linewidth}{@{}L{38mm}X@{}}
\toprule
Documento & PU-EST-001, versão 0.1 (rascunho para decisão)\\
Data & 23 de setembro de 2026\\
Autor & Lucas Salomão\\
Relaciona-se a & PU-DOC-001: \rf{RNF-01}, riscos \cd{R-01} e \cd{R-17},
  Capítulo 17 (SLOs)\\
Estado & Nenhum teste executado. Os números são estimativas com premissas
  declaradas, feitas para serem substituídas por medição.\\
\bottomrule
\end{tabularx}
\end{titlepage}

\pagenumbering{roman}
\tableofcontents
\cleardoublepage
\pagenumbering{arabic}
\setcounter{page}{1}
\pagestyle{unoomain}

% ============================================================================
\chapter{Resumo executivo}
% ============================================================================

\textbf{A pergunta}: a rede informou 800 mil usuários. Como testar a carga de
uma plataforma quando é óbvio que esses 800 mil não entram ao mesmo tempo?

\textbf{A resposta curta}: o número de cadastrados não é a carga. A carga é a
\emph{taxa de chegada} de sessões e de requisições, e ela se deriva de três
coisas: quantos estudantes estão na escola no turno mais cheio, quantas turmas
usam uma missão na mesma aula, e quanto tempo dura uma sessão. A Lei de Little
($N = \lambda W$) liga essas grandezas~\cite{little1961}, e a literatura de
teste de desempenho manda modelar \emph{sessões} em vez de \enquote{usuários
simultâneos}~\cite{meier2007cap12}, com pelo menos três cenários de uso (esperado,
melhor e pior caso).

\begin{destaque}
Com as premissas do Capítulo~\ref{cap:modelo}, a rede de 800 mil gera, no
horário mais cheio:
\begin{itemize}[nosep]
  \item \textbf{cenário esperado}: \textasciitilde\num{50000} estudantes
    simultâneos em missão, \textasciitilde210 \emph{logins} por segundo no
    primeiro minuto da aula e \textasciitilde\num{40000} requisições por segundo
    à API em regime;
  \item \textbf{cenário baixo}: metade disso; \textbf{cenário alto}: o dobro;
  \item e um \textbf{teste de ruptura} que continua subindo até quebrar, para
    saber quanta folga existe além do cenário alto.
\end{itemize}
\end{destaque}

Três conclusões mudam o plano, e aparecem antes de qualquer teste rodar:

\begin{enumerate}
  \item \textbf{A premissa de arquitetura foi escrita para outra escala.} A
    PU-DOC-001 dimensiona \textasciitilde\num{10000} estudantes e
    \num{2000}--\num{4000} simultâneos \emph{por contrato}, numa
    \enquote{instância pequena} de Cloud SQL. Uma rede de 800 mil numa única
    \emph{stack} é 80 vezes isso. Antes do teste, é preciso decidir se a rede é
    um contrato (uma \emph{stack}) ou vários (Capítulo~\ref{cap:decisoes}).
  \item \textbf{O código já aponta os gargalos prováveis}, e o teste deve
    começar por eles: o \emph{hash} de senha roda em JavaScript no
    \emph{event loop} (\textasciitilde100\,ms por \emph{login}); o
    \emph{polling} de 5\,s por estudante passa pela guarda de liberação, que
    faz de 5 a 6 \emph{queries}; o painel do professor relê todos os eventos da
    atividade a cada 5\,s; e o \emph{pool} de conexões do Prisma está no padrão
    (Capítulo~\ref{cap:hipoteses}).
  \item \textbf{O gargalo do estudante pode nem estar na nuvem.} A média de
    banda nas escolas públicas monitoradas é de 0,26\,Mbps por aluno, e só 11\%
    delas atingem a meta de 1\,Mbps por aluno~\cite{nicbr2024}. O teste de
    navegador precisa simular a rede da escola, não a fibra do escritório.
\end{enumerate}

\textbf{O que se propõe}: k6 como ferramenta, carga em \emph{modelo aberto}
(taxa de chegada, não número fixo de usuários virtuais), sete testes em
sequência (\emph{smoke}, capacidade unitária, carga média, a \enquote{parede
das 8h}, estresse e ruptura, resistência e IA isolada), mais uma pequena
fração de navegadores reais medindo Web Vitals. Os critérios de aprovação são
os SLOs que a plataforma já declarou no Capítulo 17 da PU-DOC-001.

% ============================================================================
\chapter{Por que 800 mil não é o número do teste}
\label{cap:porque}
% ============================================================================

\section{Cadastrado, ativo, simultâneo}

São três populações diferentes, e confundir as duas primeiras com a terceira é
o erro mais comum ao planejar um teste de carga:

\begin{tabularx}{\linewidth}{@{}L{32mm}X L{30mm}@{}}
\toprule
\headrow \thd{População} & \thd{Definição} & \thd{Ordem de grandeza}\\
\midrule
Cadastrados & Contas existentes no banco. Determinam o \emph{volume de dados}
  (tamanho de tabela, índice, \emph{backup}), não a carga. & \num{800000}\\
\zebra Ativos no dia & Quem entra pelo menos uma vez num dia letivo. Depende da
  adoção pedagógica: quantas aulas por semana usam uma missão. & dezenas a
  centenas de milhares\\
Simultâneos & Quem tem uma sessão aberta \emph{no mesmo instante}. É o que
  ocupa instância, conexão e memória. & dezenas de milhares\\
\bottomrule
\end{tabularx}

O guia de testes de desempenho da Microsoft é direto sobre isso: em vez de
modelar a experiência com base em usuários simultâneos, \enquote{\emph{it is
more useful to base your model on user sessions}}, e o exemplo do capítulo mostra
23 sessões numa hora produzindo de 9 a 10 usuários simultâneos~\cite{meier2007cap12}.
O \emph{syllabus} de desempenho do ISTQB faz o mesmo alerta com um exemplo
numérico: 500 usuários que fazem uma consulta curta por minuto geram
\num{30000} consultas por hora; os mesmos 500 fazendo uma por hora geram 500 ---
uma diferença de 60 vezes com o \enquote{mesmo número de usuários}~\cite{istqb2018pt}.

\section{A Lei de Little}

A ponte entre sessões e simultaneidade é a Lei de Little~\cite{little1961}:
\begin{equation}
  N = \lambda \cdot W
\end{equation}
em que $N$ é o número médio de sessões abertas, $\lambda$ é a taxa de chegada
de sessões (sessões por segundo) e $W$ é a duração média de uma sessão. Ela vale
para qualquer sistema estável, sem hipótese sobre a distribuição das chegadas.
O ISTQB traz a mesma relação escrita para usuários virtuais: vazão $=$
usuários $/$ (tempo de processamento $+$ tempo de pensamento)~\cite{istqb2018pt}.

No contexto da plataforma, isso se lê assim: se as turmas abrem missões de 45
minutos, uma taxa de chegada de 18 estudantes por segundo sustenta cerca de
\num{50000} sessões simultâneas ($18 \times 2700\,\text{s} \approx \num{48600}$).
E o inverso: se a concentração do início da aula faz $\lambda$ disparar por
alguns minutos, é esse disparo --- não a média --- que derruba o \emph{login}.

\section{Modelo aberto, não fechado}

Há duas formas de gerar carga~\cite{schroeder2006}:
\begin{itemize}
  \item \textbf{modelo fechado}: um número fixo de usuários virtuais, cada um
    só manda a próxima requisição quando a anterior volta;
  \item \textbf{modelo aberto}: as chegadas acontecem numa taxa definida,
    \emph{independentemente} de o servidor ter respondido.
\end{itemize}

Schroeder, Wierman e Harchol-Balter mostram que, para a mesma carga, o tempo de
resposta médio é \emph{significativamente menor} no modelo fechado do que no
aberto~\cite{schroeder2006}. O motivo é que, no modelo fechado, um servidor
lento \emph{desacelera o próprio gerador}: os usuários virtuais esperam, a taxa
de chegada cai, e a latência medida sai otimista. Gil Tene chamou esse efeito de
\emph{coordinated omission}: a maior parte dos artefatos de alta latência do
servidor simplesmente não é registrada~\cite{tene2013,tenewrk2}. A documentação
do k6 descreve o mesmo sintoma e recomenda os executores de taxa de chegada
(\cd{constant-arrival-rate} e \cd{ramping-arrival-rate}) para o modelo
aberto~\cite{k6openclosed}.

\begin{decisao}
\textbf{O teste usa o modelo aberto.} As crianças não esperam a API responder
para chegar à escola: às 7h30 elas entram na sala e o professor pede que abram
a missão, com o servidor lento ou não. Isso é um sistema aberto por natureza. O
modelo fechado entra só onde ele descreve o comportamento real --- o
\emph{polling} de uma aba já aberta, que espera a resposta anterior.
\end{decisao}

% ============================================================================
\chapter{Da rede ao número: o modelo de carga}
\label{cap:modelo}
% ============================================================================

\section{Premissas}

Toda estimativa abaixo é uma cadeia de premissas. Elas estão numeradas para que
a rede possa corrigir cada uma isoladamente, e o resultado se recalcula sem
reescrever o método.

\begin{tabularx}{\linewidth}{@{}L{9mm}L{50mm}L{22mm}X@{}}
\toprule
\headrow \thd{ID} & \thd{Premissa} & \thd{Valor} & \thd{Origem}\\
\midrule
\endhead
P1 & Usuários cadastrados & \num{800000} & Informado pela rede.\\
\zebra P2 & Fração de estudantes & 95\% (\num{760000}) & Premissa; o resto são
  professores e gestores. \textbf{Pedir o número à rede.}\\
P3 & Estudantes no turno mais cheio & 55\% (\textasciitilde\num{420000}) &
  Premissa. O Censo Escolar 2024 mostra que 99,9\% dos anos finais do
  fundamental e 82,5\% do ensino médio estudam no diurno, e que 22,9\% das
  matrículas públicas são de tempo integral~\cite{inep2024}; a divisão
  manhã/tarde não é publicada em agregado e tem de vir dos microdados ou da
  rede.\\
\zebra P4 & Tamanho médio de turma & 30 & Premissa, coerente com o \cd{R-17}
  da PU-DOC-001.\\
P5 & Aulas por turno e por semana & 5 por dia, 25 por semana & Premissa de
  grade típica.\\
\zebra P6 & Aulas por semana, por turma, que usam uma missão ($k$) & 1, 2 ou 4
  & \textbf{É a premissa que mais pesa}, e é pedagógica: depende de quantas
  missões o currículo da rede encaixa por semana.\\
P7 & Fator de concentração ($f$) & 1,5 & Premissa: a aula mais disputada tem
  50\% mais turmas usando a plataforma que a média (primeira aula, por
  exemplo).\\
\zebra P8 & Duração média de uma sessão ($W$) & 45 min & Uma hora-aula de 50
  minutos, menos a acomodação da turma.\\
P9 & Janela de entrada & 50\% nos 2 primeiros minutos, o resto em 8 & Premissa
  do \enquote{podem abrir agora} do professor. É a \enquote{parede} descrita na
  PU-DOC-001.\\
\bottomrule
\end{tabularx}

\section{Simultâneos no horário mais cheio}

A fração de turmas numa aula com missão é $k/25$, e a simultaneidade no horário
mais cheio é:
\begin{equation}
  N_{\text{sim}} = E_{\text{turno}} \times \frac{k}{25} \times f
\end{equation}

\begin{tabularx}{\linewidth}{@{}L{34mm}C{16mm}C{26mm}C{24mm}X@{}}
\toprule
\headrow \thd{Cenário} & \thd{$k$} & \thd{Simultâneos} & \thd{Turmas ativas} &
  \thd{Leitura}\\
\midrule
A --- baixo & 1 & \num{25200} & 840 & Início do contrato, adoção parcial.\\
\zebra B --- esperado & 2 & \num{50400} & \num{1680} & Duas missões por turma
  por semana.\\
C --- alto & 4 & \num{100800} & \num{3360} & Adoção plena, ou um
  \enquote{dia D} pedido pela secretaria.\\
\zebra D --- referência documental & --- & \num{152000} & \num{5070} & A razão
  da PU-DOC-001 (20\% dos cadastrados, o piso de
  \num{2000}--\num{4000} em \num{10000}) aplicada a \num{760000}.\\
\bottomrule
\end{tabularx}

\begin{nota}
\textbf{Por que o cenário D não é o alvo, mas entra no teste.} A razão de
20--40\% de simultâneos sobre cadastrados foi escrita para uma rede pequena e
não é plausível em 800 mil: o TIC Educação e o NIC.br mostram que só 29\% das
escolas públicas têm dispositivo para os alunos, cerca de um para cada dez, e
que o uso do celular pessoal na escola caiu de 55\% para
45\%~\cite{nicbr2024,cetic2025}. Os dispositivos são um teto natural da
simultaneidade. Mesmo assim, o cenário D serve de marco no teste de ruptura:
saber se a plataforma chega nele, ou onde para antes, custa só mais alguns
minutos de rampa.
\end{nota}

\section{A parede das 8h: taxa de \emph{login}}

Pela premissa P9, metade das sessões da aula chega nos dois primeiros minutos:

\begin{tabularx}{\linewidth}{@{}L{34mm}C{30mm}C{30mm}X@{}}
\toprule
\headrow \thd{Cenário} & \thd{Pico (2 min)} & \thd{Restante (8 min)} &
  \thd{CPU só de \emph{hash}}\\
\midrule
A & 105 \emph{logins}/s & 26/s & \textasciitilde10,5 vCPU no pico\\
\zebra B & 210 \emph{logins}/s & 52/s & \textasciitilde21 vCPU no pico\\
C & 420 \emph{logins}/s & 105/s & \textasciitilde42 vCPU no pico\\
\bottomrule
\end{tabularx}

A última coluna usa o custo medido em desenvolvimento: \textasciitilde100\,ms de
CPU por comparação \cd{bcryptjs} com custo 10
(\arq{apps/hub/lib/auth.ts:41-43}). É o pior caso --- quem já tem sessão válida
de 2\,h não refaz \emph{login} ---, e é por isso que ele é o pior caso certo
para testar: o primeiro turno do dia começa com todas as sessões vencidas.

\begin{figure}[H]
\centering
\begin{tikzpicture}
\begin{axis}[
  width=0.92\linewidth, height=58mm,
  xlabel={minutos desde o início da aula}, ylabel={\emph{logins} por segundo},
  xmin=0, xmax=14, ymin=0, ymax=240,
  grid=major, grid style={unooGrid},
  axis line style={unooAxis}, tick label style={font=\footnotesize},
  label style={font=\footnotesize},
  legend style={font=\footnotesize, draw=none, at={(0.98,0.95)}, anchor=north east}
]
\addplot[very thick, color=unooNavy, const plot]
  coordinates {(0,0) (1,210) (3,52) (11,0) (14,0)};
\addlegendentry{cenário B: chegada modelada}
\addplot[thick, dashed, color=unooLavender]
  coordinates {(0,5) (14,5)};
\addlegendentry{regime: renovação e retardatários}
\end{axis}
\end{tikzpicture}
\caption[Perfil de chegada da parede das 8h]{O perfil de chegada que o teste da
parede reproduz: um degrau de dois minutos, a cauda de oito, e o regime depois.
O Google recomenda testar tanto a rampa quanto o degrau, porque o efeito de
\emph{cache} faz os dois darem resultados diferentes~\cite{srecascading}.}
\label{fig:parede}
\end{figure}

\section{Requisições por segundo em regime}

Depois que a missão abre, cada aba de estudante gera tráfego próprio. O
levantamento do código (Capítulo~\ref{cap:hipoteses}) dá a composição:

\begin{tabularx}{\linewidth}{@{}L{46mm}L{24mm}X@{}}
\toprule
\headrow \thd{Fonte} & \thd{Taxa por aba} & \thd{Onde}\\
\midrule
\emph{Polling} de estado (fase, relógio, portão) & 1 a cada 5\,s &
  \arq{exp08ef1b}, \arq{exp02em1b}, \arq{exp06ef2b}\\
\zebra Eventos do motor (Instanooria) & lote a cada 4\,s &
  \arq{exp08ef1b/\B{}src/\B{}lib/\B{}eventos.ts:\B{}51-52}\\
Telemetria & até 6 por minuto & \arq{packages/\B{}sdk/\B{}src/\B{}telemetry.ts}\\
\zebra Ações da missão & \textasciitilde1 por minuto & rotas \cd{POST} de cada
  etapa\\
\emph{Preflight} CORS & até $\times 2$ & \arq{apps/api/src/main.ts:48}: sem
  \cd{maxAge}, o navegador guarda o \emph{preflight} por poucos segundos\\
\bottomrule
\end{tabularx}

Somado, dá de 0,35 a 1,2 requisições por segundo por aba; o valor central de
0,8 dá \textasciitilde\num{20000} req/s no cenário A,
\textasciitilde\num{40000} no B e \textasciitilde\num{80000} no C. Esse é o
número que dimensiona a API, e ele é dominado pelo \emph{polling}, não pelo que
a criança faz. A documentação do k6 recomenda justamente descrever carga de API
por taxa de requisições, e converter para taxa de iterações dividindo pelo
número de requisições por iteração~\cite{k6apitesting}.

\begin{atencao}
Estes números são ordem de grandeza. O objetivo deles é escolher a escala do
gerador e da cota de nuvem, não prever o resultado. O primeiro teste
(capacidade unitária, Seção~\ref{sec:unitaria}) substitui a estimativa pela
medida.
\end{atencao}

% ============================================================================
\chapter{O que testar}
\label{cap:oque}
% ============================================================================

\section{As jornadas}

O teste não exercita rotas soltas: exercita as jornadas que acontecem numa
aula, na proporção em que acontecem. Uma jornada é uma sequência de requisições
com tempo de pensamento entre elas~\cite{meier2007cap12}.

\begin{tabularx}{\linewidth}{@{}L{8mm}L{40mm}X L{14mm}@{}}
\toprule
\headrow \thd{J} & \thd{Jornada} & \thd{Sequência (requisições reais)} &
  \thd{Peso}\\
\midrule
\endhead
J1 & \emph{Login} do estudante &
  \rota{GET /login} \seta{} \rota{/api/auth/providers} \seta{}
  \rota{/api/auth/csrf} \seta{} \rota{POST /api/auth/callback/credentials}
  (que chama \rota{POST /auth/login} na API) \seta{} \rota{/api/auth/session}
  \seta{} \rota{GET /estudante} (que chama \rota{GET /salas/minhas}). Cerca de
  8 requisições ao hub e 2 à API. & 1 por sessão\\
\zebra J2 & Abrir a missão &
  \rota{GET /exp/:expId}: \rota{/salas/minhas}, \rota{/experiencias/:expId},
  \rota{/recursos/efetivos} (3 chamadas à API, 2 assinaturas RS256) \seta{}
  \emph{bundle} da CDN \seta{} \emph{handshake} \cd{exp:ready}/\cd{hub:init}
  \seta{} rota de entrada da experiência (passa por
  \cd{exigirMissaoLiberada}). & 1 por sessão\\
J3 & Missão sem IA (Rastro do Boato, \arq{exp06ef2b}) &
  \rota{GET /rastro-boato/sessao}, \rota{POST /entrar}, \emph{polling} de
  \rota{/gate} a cada 5\,s, \rota{POST} de auditoria, decisão, modelagem e
  reflexão. & 50\%\\
\zebra J4 & Missão com motor de eventos (Instanooria, \arq{exp08ef1b}) &
  \rota{GET /runs/boot}, \rota{POST /runs}, \rota{/entrar}, \rota{/feed},
  \emph{polling} de \rota{/fase} a cada 5\,s, \rota{POST /runs/:id/eventos} a
  cada 4\,s, \rota{student-post} (IA opcional). & 35\%\\
J5 & Missão com IA síncrona (unooGPT, \arq{exp02em1b}) &
  \rota{GET /unoogpt/sessao}, \emph{polling} de \rota{/relogio}, \rota{POST
  /unoogpt/pergunta} (10--30\,s, até 6 voltas de ferramenta). & 15\%\\
\zebra J6 & Painel do professor & Uma aba por turma ativa: monitor a cada 5\,s
  (\rota{GET /atividades/:id/monitor}) ou \cd{router.refresh()} a cada 5--10\,s,
  conforme a missão. & 1 por turma\\
J7 & Renovação do \emph{token} & \rota{GET /api/exp-token} a cada 10 minutos por
  aba aberta. & contínuo\\
\zebra J8 & Telemetria & \rota{POST /telemetry/batch} \seta{} Pub/Sub \seta{}
  \rota{POST /telemetry/ingest}. Mede-se a fila, não só a rota. & contínuo\\
\bottomrule
\end{tabularx}

Os pesos de J3 a J5 são premissa e devem refletir o calendário pedagógico da
rede no mês do teste. O que importa é que as três classes de missão estejam
presentes: a que só lê e grava, a que grava eventos em alta frequência e a que
segura a requisição por dezenas de segundos esperando a IA.

\section{Fora do escopo desta rodada}

\begin{itemize}
  \item \textbf{O provedor de identidade (WSO2).} Hoje o hub é o emissor, e o
    WSO2 entra depois. Quando entrar, precisa de teste próprio: em agosto de
    2020, o primeiro dia de aulas de distritos de San Antonio (EUA) caiu às
    6h45 porque o portal de \emph{login} não aguentou o volume, e a falha
    arrastou em cascata o Google Classroom e o Zoom, que dependiam
    dele~\cite{carnett2020}.
  \item \textbf{A CDN em si.} Ela é um serviço gerenciado com capacidade muito
    acima da nossa; o que interessa medir é a taxa de acerto de
    \emph{cache}, porque o Google avisa que nenhuma requisição tem garantia de
    sair do \emph{cache}~\cite{gcdnoverview}. Seguindo a recomendação do k6, os
    testes de protocolo excluem os \emph{assets} da CDN e os testes de
    navegador os incluem~\cite{k6websites}.
  \item \textbf{Telas de gestão} (importação, relatórios da diretoria): baixo
    volume, com limite próprio; entram só no teste de resistência como ruído de
    fundo.
\end{itemize}

% ============================================================================
\chapter{Hipóteses de gargalo que o código já aponta}
\label{cap:hipoteses}
% ============================================================================

Um teste de carga bem feito começa com hipóteses: sem elas, a equipe olha o
painel no dia e não sabe o que procurar. O levantamento do código em 23 de
setembro de 2026 aponta oito, e cada teste do Capítulo~\ref{cap:tipos} existe
para confirmar ou descartar pelo menos uma delas. \textbf{Nenhuma foi medida
ainda}: são leituras do código.

\begin{tabularx}{\linewidth}{@{}L{8mm}L{44mm}X L{22mm}@{}}
\toprule
\headrow \thd{H} & \thd{Hipótese} & \thd{Por que o código sugere} &
  \thd{Teste que decide}\\
\midrule
\endhead
H1 & O \emph{login} satura a CPU da API &
  \cd{bcryptjs} (JavaScript puro) com custo 10, no \emph{event loop}:
  \textasciitilde100\,ms por comparação, e uma instância Node faz no máximo
  \textasciitilde10 por segundo enquanto bloqueia todo o resto
  (\arq{apps/\B{}api/\B{}src/\B{}auth/\B{}auth.controller.ts:\B{}141,154-167}). &
  Parede das 8h\\
\zebra H2 & A busca do \emph{login} não tem índice &
  \cd{user.findMany(\{username, ativo\})}, mas o único índice é o
  \cd{UNIQUE (escola\_id, username)} (\arq{schema.prisma:442}). Com 800 mil
  linhas, a busca sem \cd{escola\_id} pode varrer a tabela. &
  Capacidade unitária\\
H3 & O \emph{pool} de conexões esgota antes da CPU &
  Nenhum \cd{connection\_limit} no código: vale o padrão do Prisma
  ($2 \times \text{vCPU} + 1$). Cada \cd{comTenancy} abre uma transação
  interativa (há 239 chamadas), com \cd{maxWait} de 2\,s. O
  \emph{pooler} está declarado na documentação e não aparece no repositório. &
  Carga média, estresse\\
\zebra H4 & O \emph{polling} multiplica a guarda de liberação &
  Toda rota de estudante, inclusive o \emph{polling} de 5\,s do portão, passa
  por \cd{exigirMissaoLiberada}, que faz de 5 a 6 \emph{queries}
  (\arq{apps/\B{}api/\B{}src/\B{}comum/\B{}liberacao.ts:\B{}104-130}). No cenário B, se todos
  estivessem nessa missão, seriam \textasciitilde\num{60000} \emph{queries}/s
  só de guarda. & Carga média\\
H5 & O monitor do professor cresce durante a aula &
  \rota{GET /atividades/:id/monitor} a cada 5\,s lê todas as \emph{runs} e todos
  os eventos da atividade, sem \cd{take}
  (\arq{apps/\B{}api/\B{}src/\B{}analise/\B{}metricas.ts:\B{}123-160}). O custo por chamada
  aumenta a cada minuto de aula. & Resistência\\
\zebra H6 & A IA prende a capacidade da API &
  \rota{POST /unoogpt/pergunta} leva 10--30\,s. Pela Lei de Little, 5 mil
  estudantes perguntando uma vez a cada 5 minutos, com 20\,s de espera, mantêm
  \textasciitilde330 requisições em voo; a largada simultânea da turma
  multiplica isso por 5. Cada uma ocupa um dos 80 lugares de concorrência da
  instância. & IA isolada\\
H7 & Limites em memória viram $N$ vezes o declarado &
  O \emph{rate limit} e as quotas locais são um \cd{Map} por instância
  (\arq{apps/\B{}api/\B{}src/\B{}comum/\B{}rate-limit.ts}). Com 50 instâncias, o \enquote{8
  por minuto} vira até 400. Não é gargalo: é proteção que some sob carga. &
  Estresse\\
\zebra H8 & O hub assina demais &
  Cada chamada do hub à API assina um JWT RS256 novo
  (\arq{apps/\B{}hub/\B{}lib/\B{}api-servidor.ts:\B{}80,121}), e cada abertura de missão faz
  duas. A API verifica cada uma. & Parede das 8h\\
\bottomrule
\end{tabularx}

\begin{nota}
O teste não corrige nada disso: ele mede. A regra do projeto é \enquote{nenhuma
peça de infraestrutura sem gargalo medido} (RE-01 da PU-DOC-001), e esta lista
existe para que o gargalo, quando aparecer, já tenha um nome e um arquivo.
\end{nota}

% ============================================================================
\chapter{Os tipos de teste e a ordem de execução}
\label{cap:tipos}
% ============================================================================

A classificação segue a documentação do k6~\cite{k6testtypes} e o
\emph{syllabus} do ISTQB~\cite{istqb2018pt}. A ordem importa: cada teste só roda
se o anterior passou, porque um teste de estresse num sistema que falha no
\emph{smoke} só mede o defeito já conhecido.

\begin{tabularx}{\linewidth}{@{}L{6mm}L{30mm}L{30mm}X@{}}
\toprule
\headrow \thd{\#} & \thd{Teste} & \thd{Carga e duração} & \thd{Pergunta que
  responde}\\
\midrule
\endhead
1 & \emph{Smoke} & 5--10 usuários, 2 min & O roteiro funciona? \emph{Login},
  \emph{token}, missão e painel respondem 2xx?\\
\zebra 2 & Capacidade unitária & 1 instância fixa da API e do hub, rampa até
  quebrar & Quantas req/s \emph{uma} instância aguenta dentro do SLO? É o
  número que dimensiona tudo.\\
3 & Carga média & Cenário B, 30--60 min & O sistema se mantém dentro do SLO na
  carga esperada?\\
\zebra 4 & Parede das 8h (\emph{spike}) & Degrau da
  Figura~\ref{fig:parede}, 15 min & O \emph{autoscaling} acompanha o degrau? Há
  429 ou 5xx no primeiro minuto?\\
5 & Estresse e ruptura (\emph{breakpoint}) & Cenário C, depois rampa até D e
  além & Onde quebra, e como? Degrada com 429 e volta sozinho, ou cai em
  cascata?\\
\zebra 6 & Resistência (\emph{soak}) & Cenário A ou B, 5 h (um turno inteiro) &
  Vazamento de memória, \emph{pool} que não devolve conexão, fila de telemetria
  que cresce, monitor do professor ficando lento.\\
7 & IA isolada & Turmas inteiras de J5 largando juntas & \cd{R-17}: quantos
  turnos simultâneos cabem na cota do Vertex e na concorrência da API?\\
\bottomrule
\end{tabularx}

\section{Capacidade unitária: o teste que mais ensina}
\label{sec:unitaria}

Com o \emph{autoscaling} desligado (mínimo e máximo de instâncias iguais a 1),
sobe-se a taxa de chegada até que o p95 passe do SLO ou a taxa de erro passe de
1\%. O resultado é a capacidade $c$ de uma instância, em req/s. A partir dela:
\begin{equation}
  \text{instâncias necessárias} = \frac{\text{req/s do cenário}}{c \times 0{,}6}
\end{equation}
em que 0,6 vem do alvo de 60\% de CPU e de concorrência que o
\emph{autoscaler} do Cloud Run persegue~\cite{gcrautoscaling}. O Google
recomenda repetir o teste em cargas diferentes para confirmar que o serviço
escala de forma linear~\cite{gcrloadtest} --- e ele não escala: a Lei Universal
de Escalabilidade de Gunther descreve por que contenção e coerência fazem a
curva dobrar~\cite{gunther2007}. É por isso que o teste 3 existe mesmo depois
do 2.

\section{Carga híbrida: protocolo mais navegador}

O k6 recomenda gerar a maior parte da carga em nível de protocolo e rodar,
junto, um número pequeno de navegadores reais~\cite{k6websites}. Os navegadores
medem o que o protocolo não vê --- o LCP, o INP e o tempo até o
\emph{handshake} da ponte --- e o módulo de navegador do k6 exporta essas
métricas diretamente~\cite{k6browsermetrics}. Proposta: 50 a 100 navegadores
durante os testes 3 e 4, com a rede limitada a um perfil de escola (ver
Seção~\ref{sec:banda}).

% ============================================================================
\chapter{Critérios de aprovação e de parada}
\label{cap:criterios}
% ============================================================================

\section{Aprovação: os SLOs que a plataforma já declarou}

O teste não inventa critério: usa os SLOs do Capítulo 17 da PU-DOC-001, porque
são as promessas que a operação vai cobrar depois. O livro de SRE do Google
insiste em percentis, não em médias, porque a média esconde a cauda que o
usuário sente~\cite{sreslo}. No k6, cada linha vira um \emph{threshold}, que é
o critério de aprovação ou reprovação do teste~\cite{k6thresholds}.

\begin{tabularx}{\linewidth}{@{}L{44mm}L{36mm}X@{}}
\toprule
\headrow \thd{Critério} & \thd{Limite} & \thd{Origem}\\
\midrule
\endhead
Erros 5xx na API e no hub & $< 0{,}5\%$ & SLO de disponibilidade de 99,5\%\\
\zebra Latência de \rota{/runs/:id/feed}, \rota{/experiencias/:id},
  \rota{/auth/*} & p95 $< 800$\,ms & SLO de latência da API\\
\emph{Login} completo (J1 inteira) & p95 $< 3$\,s & Proposta deste estudo\\
\zebra \emph{Handshake} da ponte & 99\% em até 8\,s & SLO da ponte\\
Turno de IA & 98\% concluídos; p95 $< 30$\,s & SLO de IA e alerta B\\
\zebra Fila de telemetria & mensagem mais antiga $< 5$\,min & SLO de
  ingestão\\
LCP e INP (navegador) & p75: LCP $\leq 2{,}5$\,s, INP $\leq 200$\,ms &
  Core Web Vitals~\cite{webvitals}\\
\zebra Respostas 429 no degrau & $< 1\%$ no primeiro minuto & O Cloud Run
  segura requisições por até 10\,s sem instância livre e depois responde
  429~\cite{gcrautoscaling}\\
\bottomrule
\end{tabularx}

\section{Parada: quando interromper}

\begin{itemize}
  \item Taxa de erro acima de 5\% por mais de um minuto: o k6 aborta sozinho
    com \cd{abortOnFail}~\cite{k6thresholds}.
  \item CPU do Cloud SQL acima de 90\% por cinco minutos, ou conexões no
    limite: o banco é compartilhado por todas as escolas e não se recupera
    sozinho de uma fila de conexões.
  \item \textbf{O gerador passou de 80\% de CPU ou 90\% de memória.} Aí o
    resultado é do gerador, não do sistema~\cite{k6largetests}.
  \item Gasto de IA acima do teto aprovado para o teste.
\end{itemize}

\section{O que observar do lado do servidor}

A observabilidade do Capítulo 17 já está no ar e é o que torna o teste legível.
Durante cada rodada:

\begin{itemize}
  \item \textbf{Cloud Run}: instâncias, concorrência por instância, CPU,
    memória, latência de \emph{startup}, contagem de 429;
  \item \textbf{Cloud SQL}: CPU, conexões ativas, esperas por \emph{lock},
    \cd{pg\_stat\_statements} (as dez \emph{queries} mais caras), Query
    Insights;
  \item \textbf{API}: erros \cd{P2024} do Prisma (espera por conexão do
    \emph{pool}), \emph{queries} acima de 500\,ms (o log lento já existe),
    latência de IA por rota;
  \item \textbf{Pub/Sub}: \cd{oldest\_unacked\_message\_age} e fila morta;
  \item \textbf{Vertex}: respostas 429 por cota.
\end{itemize}

O Google recomenda exportar os logs para o BigQuery e analisar a latência por
lá depois do teste~\cite{gcrloadtest}.

% ============================================================================
\chapter{Ambiente, dados e ferramentas}
\label{cap:ambiente}
% ============================================================================

\section{Onde rodar}

\begin{decisao}
\textbf{Na \emph{stack} de produção da rede, antes do lançamento, com dados
sintéticos, e apagando tudo depois.} Até novembro de 2026 a produção não tem
nenhum estudante real: é a única janela em que dá para testar a configuração
exata que vai atender a rede (\emph{tier} do banco, cotas, rede, CDN) sem
arriscar dado de menor. Depois do lançamento, o teste passa para uma cópia da
\emph{stack} em homologação, e nunca volta para produção.
\end{decisao}

\begin{itemize}
  \item \textbf{Região}: a de produção (\cd{southamerica-east1}). O Google
    sugere \cd{europe-west1} para testes de carga no Cloud Run por causa da
    cota~\cite{gcrloadtest}, mas testar noutra região mede outro sistema; a
    alternativa é pedir aumento de cota na região real.
  \item \textbf{Cotas}: o máximo padrão do Cloud Run é de 100 instâncias por
    serviço, e o Google pede que se fale com a equipe da conta quando o teste
    for passar disso~\cite{gcrloadtest}. É provável que o cenário B já passe.
  \item \textbf{Aviso ao Google}: nenhuma fonte oficial exige aviso prévio para
    teste de carga nos próprios projetos; o FAQ de segurança dispensa até o de
    teste de invasão, desde que respeitada a política de uso
    aceitável~\cite{gcppentestfaq}. O que existe é a conversa de cota acima.
  \item \textbf{Gerador}: não usar o Pub/Sub para gerar carga, que o próprio
    Google desaconselha~\cite{gcrloadtest}. Usar máquinas na mesma região, fora
    da VPC da aplicação, para que o tráfego entre pelo \emph{load balancer} como
    o de uma escola.
\end{itemize}

\section{Dados sintéticos, nunca reais}

\begin{regra}
Nenhum dado real de estudante entra no teste, nem anonimizado. A LGPD exige que o
tratamento de dados de crianças e adolescentes seja feito no melhor interesse
deles (art. 14) e limitado ao mínimo necessário (art. 6º, III), e o dado
anonimizado volta a ser pessoal se a anonimização for revertida
(art. 12)~\cite{lgpd2018}. Um teste de carga não precisa de nenhuma pessoa real.
\end{regra}

A massa é gerada pelo mesmo caminho da produção --- o importador (ADR-54 e
ADR-59), que é quem escreve pessoa ---, com o gerador determinístico que já
existe no repositório (\arq{apps/\B{}api/\B{}prisma/\B{}pessoas-demo.ts}):

\begin{itemize}
  \item \num{800000} usuários em \textasciitilde\num{2200} escolas sintéticas
    (\textasciitilde350 estudantes por escola), turmas de 30, professores
    lotados;
  \item liberações criadas para todas as turmas que o cenário usa, porque sem
    elas a guarda responde 403 e o teste mede o 403;
  \item \textbf{um único \emph{hash} de senha reaproveitado}: 800 mil
    \cd{bcrypt} de custo 10 custariam \textasciitilde22 horas de CPU; como o
    sal viaja dentro do \emph{hash}, a mesma string vale para todos;
  \item o arquivo de credenciais do gerador fica fora do repositório.
\end{itemize}

Carregar os 800 mil pelo importador é, de quebra, o primeiro teste: é o volume
de uma importação inicial real.

\section{IA: simulada em escala, real em amostra}

A IA custa dinheiro por \emph{token} e tem cota própria. Uma partida da
unooGPT custou \num{71397} \emph{tokens} (medido em 27/08/2026); 5 mil estudantes
jogando ao mesmo tempo seriam \textasciitilde9 milhões de \emph{tokens} por
minuto. Proposta:

\begin{itemize}
  \item nos testes 3 a 6, o \emph{provider} \cd{mock} com \emph{latência
    sintética} de 10--30\,s --- o que interessa ali é quanto a espera prende a
    API, não o texto (a memória do projeto registra que o \emph{mock} atual
    lança em \cd{generateJson}, o que precisa ser resolvido antes);
  \item no teste 7, o Vertex real, com um teto de gasto fixado antes, para
    medir a cota e a latência sob concorrência. O capítulo de sobrecarga do
    livro de SRE descreve as ferramentas para quando a cota morder: cota por
    cliente, \emph{throttling} adaptativo e níveis de
    criticidade~\cite{sreoverload}.
\end{itemize}

\section{Banda da escola}
\label{sec:banda}

O MEC fixou em 2024 a meta de 1\,Mbps por aluno no turno de maior
matrícula~\cite{enec2024}; o NIC.br mediu que só 11\% das escolas monitoradas
chegam lá e que a média nacional é de 0,26\,Mbps por aluno~\cite{nicbr2024}.
Numa escola de 600 estudantes no turno, isso é \textasciitilde156\,Mbps. Se três
turmas de 30 abrem juntas uma missão de \textasciitilde8\,MB, são
\textasciitilde720\,MB, ou \textasciitilde37\,s com o \emph{link} inteiro
dedicado a isso. Os vídeos da \arq{exp04ef1b} somam 52\,MB.

Por isso os navegadores do teste híbrido rodam com a rede limitada, e o LCP
entra como critério. O Khan Academy atribuiu ter aguentado 2,5 vezes o tráfego
em duas semanas, em 2020, a exatamente duas coisas: servidor sem estado e CDN
agressiva~\cite{kosarchyn2020} --- a mesma aposta da ADR-01 da plataforma.

\section{Ferramenta}

\begin{tabularx}{\linewidth}{@{}L{22mm}L{26mm}L{38mm}X@{}}
\toprule
\headrow \thd{Ferramenta} & \thd{Roteiro em} & \thd{Modelo aberto} &
  \thd{Navegador real}\\
\midrule
\textbf{k6} & JavaScript & nativo (\cd{arrival-rate})~\cite{k6openclosed} &
  sim, com Web Vitals~\cite{k6browsermetrics}\\
\zebra Gatling & Scala, Java, Kotlin & \cd{injectOpen}
  explícito~\cite{gatlinginjection} & não\\
Locust & Python & por taxa de usuários~\cite{locustdocs} & não\\
\zebra JMeter & GUI, Java & por plano & não; trabalha em nível de protocolo e
  não executa JavaScript~\cite{jmeter}\\
\bottomrule
\end{tabularx}

\begin{decisao}
\textbf{k6.} O roteiro é JavaScript, a linguagem do projeto inteiro; o modelo
aberto é nativo; o mesmo roteiro roda os navegadores; e a PU-DOC-001 já o
nomeia no \cd{R-01}. A execução distribuída pode ser feita com várias máquinas
rodando k6 por fatia do cenário, ou com um serviço gerenciado. O tutorial do
Google de teste distribuído com Locust em GKE foi arquivado em agosto de 2026 e
não deve ser seguido~\cite{gcpgkelocust}.
\end{decisao}

\textbf{Escala do gerador.} No modelo aberto, o número de usuários virtuais que
o k6 precisa é a taxa de chegada vezes a duração da iteração (Little de novo).
Modelar a sessão inteira de 45 minutos como uma iteração exigiria
\textasciitilde\num{50000} VUs no cenário B; a 1--5\,MB por VU, isso é de
50 a 250\,GB de memória~\cite{k6largetests}. A saída prática é separar os
cenários: a jornada completa (J1 e J2) como iteração curta, e o regime (o
\emph{polling}, os eventos, o painel) como taxa de requisições por rota, na
proporção medida. O Apêndice~\ref{ap:k6} mostra o esqueleto.

% ============================================================================
\chapter{Contas de capacidade para levar ao teste}
\label{cap:contas}
% ============================================================================

\section{Conexões com o banco}

Duas regras do Google limitam as conexões: cada instância do Cloud Run abre no
máximo 100 conexões com o Cloud SQL~\cite{gcsqlconnectrun}, e o
\cd{max\_connections} padrão do Postgres depende da memória da instância ---
100 entre 3,75 e 6\,GB, 400 entre 7,5 e 15\,GB, 500 entre 15 e 30\,GB,
\num{1000} a partir de 120\,GB~\cite{gcsqlflags}. A condição a respeitar é:
\begin{equation}
  \text{máx. instâncias da API} \times \text{\emph{pool} por instância}
  \;\leq\; \texttt{max\_connections} - \text{reserva}
\end{equation}
Com o \emph{pool} padrão do Prisma numa instância de 2 vCPU (5 conexões), 100
instâncias pedem 500 conexões --- mais que qualquer \emph{tier} até 30\,GB. Ou
o \emph{pooler} existe de fato (a documentação diz que sim, o repositório não
mostra), ou o máximo de instâncias fica preso pelo banco. O Google recomenda
\emph{pool} pequeno por instância, com \emph{timeout} e \emph{backoff}
exponencial~\cite{gcsqlmanageconn}. Este é o primeiro número a confirmar, antes
do primeiro teste com mais de uma instância.

\section{IA pela Lei de Little}

Requisições de IA em voo $=$ taxa de perguntas $\times$ duração da pergunta.
Com 5 mil estudantes na unooGPT (15\% do cenário B, arredondado), uma pergunta a
cada 5 minutos por estudante e 20\,s de espera:
\begin{equation}
  N_{\text{IA}} = \frac{5000}{300\,\text{s}} \times 20\,\text{s} \approx 333
\end{equation}
Na largada, quando 30 estudantes de cada turma perguntam no mesmo minuto, o
número sobe para \textasciitilde\num{1700}. Com concorrência 80 por instância,
são de 4 a 21 instâncias fazendo nada além de esperar o Vertex --- e cada
uma dessas requisições é também um lugar que o estudante do lado, fazendo uma
chamada de 50\,ms, não consegue ocupar. É o \cd{R-17}; a saída já registrada é
\emph{streaming} no \emph{provider}.

\section{\emph{Login}}

A capacidade de \emph{login} por instância é limitada por
$1 / 0{,}1\,\text{s} = 10$ por vCPU enquanto o \emph{hash} rodar no
\emph{event loop}. Para os 210/s do cenário B, são pelo menos 21 vCPU só de
\emph{hash}, e cada \emph{login} bloqueia por 100\,ms todas as outras
requisições da mesma instância. Se H1 se confirmar, é a primeira correção a
fazer, e ela está dentro da API.

% ============================================================================
\chapter{Decisões que o teste não toma}
\label{cap:decisoes}
% ============================================================================

O teste mede. Algumas perguntas vêm antes dele e não são técnicas:

\begin{enumerate}
  \item \textbf{A rede de 800 mil é um contrato ou vários?} A ADR-04 define
    \enquote{1 contrato $=$ 1 \emph{stack}}, com a escala pensada para
    \textasciitilde\num{10000} estudantes. Uma \emph{stack} para 800 mil muda o
    \emph{tier} do banco, as cotas e o custo, e concentra o risco: um incidente
    para a rede inteira. Dividir por regional (várias \emph{stacks} para o mesmo
    cliente) mantém cada uma perto da escala desenhada, e custa operar
    várias. \textbf{O teste deve rodar sobre a forma escolhida.}
  \item \textbf{Qual é a adoção esperada ($k$)?} Uma, duas ou quatro missões por
    turma por semana mudam a carga em quatro vezes. Isso é do planejamento
    pedagógico da rede.
  \item \textbf{Qual é o horário de início das aulas e a divisão entre
    turnos?} As fontes nacionais não publicam isso em agregado; a rede tem o
    dado.
  \item \textbf{Quanto se pode gastar no teste?} Instâncias, banco maior
    durante o teste e \emph{tokens} de IA do teste 7 precisam de teto aprovado.
\end{enumerate}

\begin{nota}
Ao responder, atualize as premissas P2, P3, P6 e P7 do
Capítulo~\ref{cap:modelo}: as tabelas de cenário se recalculam com a mesma
fórmula.
\end{nota}

% ============================================================================
\chapter{Plano de execução}
\label{cap:plano}
% ============================================================================

Com produção prevista para novembro de 2026, sobram cerca de seis semanas.

\begin{tabularx}{\linewidth}{@{}L{20mm}X L{36mm}@{}}
\toprule
\headrow \thd{Semana} & \thd{Entrega} & \thd{Sai daqui}\\
\midrule
\endhead
1 (até 02/10) & Respostas do Capítulo~\ref{cap:decisoes}; conferir
  \emph{pooler}, \cd{connection\_limit} e \emph{tier}; pedido de cota; massa
  sintética carregada pelo importador; \emph{mock} de IA com latência. &
  Ambiente pronto\\
\zebra 2 & Roteiros k6 das jornadas J1--J8; teste 1 (\emph{smoke}); teste 2
  (capacidade unitária). & Capacidade $c$ por instância\\
3 & Correções das hipóteses confirmadas no teste 2 (provavelmente H1, H2,
  H3); repetir o teste 2. & $c$ corrigido\\
\zebra 4 & Teste 3 (carga média) e teste 4 (parede das 8h), com navegadores. &
  Aprovação no cenário B\\
5 & Teste 5 (estresse e ruptura) e teste 6 (resistência de 5\,h). & Ponto de
  ruptura; vazamentos\\
\zebra 6 & Teste 7 (IA real); relatório final; atualização da PU-DOC-001
  (\rf{RNF-01}, \cd{R-01}, \cd{R-17}, Capítulo 6); limpeza da massa
  sintética. & \cd{R-01} fechado ou reclassificado\\
\bottomrule
\end{tabularx}

\begin{atencao}
A semana 3 é a que costuma faltar. Todo teste de carga que acha um gargalo
precisa de tempo para corrigir e \emph{repetir}; um cronograma sem essa semana
só descobre o problema, não o resolve. Os incidentes recorrentes no primeiro dia
de inscrição do Enem, atribuídos pela imprensa ao volume de acessos
simultâneos, mostram que o risco se repete quando nada muda entre uma edição e
a próxima~\cite{techtudo2024,querobolsa2026}.
\end{atencao}

% ============================================================================
\appendix
\chapter{Esqueleto do roteiro k6}
\label{ap:k6}
% ============================================================================

Esqueleto para orientar a escrita, \textbf{não testado}. Os nomes de rota são
os do código em 23 de setembro de 2026, e o fluxo de \emph{login} segue o
NextAuth com o \emph{provider} de credenciais.

\begin{lstlisting}[language=tsunoo]
import http from "k6/http";
import { check, sleep } from "k6";
import { SharedArray } from "k6/data";

const HUB = __ENV.HUB_URL;
const API = __ENV.API_URL;
const contas = new SharedArray("contas", () =>
  open("./contas-sinteticas.csv").split("\n").slice(1).map((l) => l.split(",")));

export const options = {
  scenarios: {
    // J1 + J2: a parede das 8h (cenario B), modelo aberto.
    parede: {
      executor: "ramping-arrival-rate",
      exec: "entrarNaAula",
      startRate: 0, timeUnit: "1s",
      preAllocatedVUs: 2000, maxVUs: 8000,
      stages: [
        { target: 210, duration: "1m" },  // o degrau
        { target: 210, duration: "1m" },
        { target: 52,  duration: "30s" },
        { target: 52,  duration: "7m30s" },
        { target: 0,   duration: "1m" },
      ],
    },
    // Regime: polling do portao, na taxa por rota (e nao por usuario).
    regime: {
      executor: "constant-arrival-rate",
      exec: "pollingDoPortao",
      rate: 10000, timeUnit: "1s", duration: "30m",
      preAllocatedVUs: 3000, maxVUs: 12000,
      startTime: "10m",
    },
  },
  thresholds: {
    http_req_failed: [{ threshold: "rate<0.005", abortOnFail: true,
                        delayAbortEval: "1m" }],
    "http_req_duration{grupo:api-slo}": ["p(95)<800"],
    "http_req_duration{grupo:login}":   ["p(95)<3000"],
    dropped_iterations: ["count<100"],  // o gerador nao acompanhou
  },
};

export function entrarNaAula() {
  const [usuario, senha, expId] = contas[Math.floor(Math.random() * contas.length)];
  const csrf = http.get(`${HUB}/api/auth/csrf`, { tags: { grupo: "login" } })
                   .json("csrfToken");
  const r = http.post(`${HUB}/api/auth/callback/credentials`,
    { csrfToken: csrf, username: usuario, password: senha },
    { tags: { grupo: "login" }, redirects: 0 });
  check(r, { "login aceito": (x) => x.status === 302 || x.status === 200 });
  http.get(`${HUB}/estudante`, { tags: { grupo: "login" } });
  http.get(`${HUB}/exp/${expId}`, { tags: { grupo: "abrir" } });
  const token = http.get(`${HUB}/api/exp-token`).json("token");
  http.get(`${API}/rastro-boato/sessao`, {
    headers: { Authorization: `Bearer ${token}` }, tags: { grupo: "api-slo" } });
}

export function pollingDoPortao() {
  // Tokens pre-emitidos para evitar que o regime refaca login.
  // ...
}
\end{lstlisting}

O detalhe que mais importa está em \cd{dropped\_iterations}: no modelo aberto,
quando os usuários virtuais pré-alocados acabam, o k6 descarta iterações em vez
de desacelerar. Contá-las é o que separa \enquote{o sistema aguentou} de
\enquote{o gerador não conseguiu mandar}.

% ============================================================================
\chapter{Referências de método sem citação direta no texto}
% ============================================================================

As três obras abaixo são a base clássica de caracterização de carga e
planejamento de capacidade, e sustentam o método do Capítulo~\ref{cap:modelo}
sem que uma passagem específica tenha sido citada: \citeonline{jain1991},
\citeonline{menasce2002} e a norma de qualidade de produto
\citeonline{iso25010}, cuja característica de \emph{eficiência de
desempenho} (comportamento no tempo, uso de recursos e capacidade) é a que este
teste verifica. O modelo RAIL~\cite{rail} complementa os Web Vitals com os
limites de resposta de interação. A classificação do \emph{syllabus} ISTQB e a
do guia da Microsoft~\cite{meier2007} coincidem com a do k6 usada no
Capítulo~\ref{cap:tipos}.

\bibliography{referencias}

\end{document}
